% Options for packages loaded elsewhere
\PassOptionsToPackage{unicode}{hyperref}
\PassOptionsToPackage{hyphens}{url}
\PassOptionsToPackage{dvipsnames,svgnames,x11names}{xcolor}
%
\documentclass[
  11pt,
]{article}
\usepackage{amsmath,amssymb}
\usepackage{iftex}
\ifPDFTeX
  \usepackage[T1]{fontenc}
  \usepackage[utf8]{inputenc}
  \usepackage{textcomp} % provide euro and other symbols
\else % if luatex or xetex
  \usepackage{unicode-math} % this also loads fontspec
  \defaultfontfeatures{Scale=MatchLowercase}
  \defaultfontfeatures[\rmfamily]{Ligatures=TeX,Scale=1}
\fi
\usepackage{lmodern}
\ifPDFTeX\else
  % xetex/luatex font selection
\fi
% Use upquote if available, for straight quotes in verbatim environments
\IfFileExists{upquote.sty}{\usepackage{upquote}}{}
\IfFileExists{microtype.sty}{% use microtype if available
  \usepackage[]{microtype}
  \UseMicrotypeSet[protrusion]{basicmath} % disable protrusion for tt fonts
}{}
\makeatletter
\@ifundefined{KOMAClassName}{% if non-KOMA class
  \IfFileExists{parskip.sty}{%
    \usepackage{parskip}
  }{% else
    \setlength{\parindent}{0pt}
    \setlength{\parskip}{6pt plus 2pt minus 1pt}}
}{% if KOMA class
  \KOMAoptions{parskip=half}}
\makeatother
\usepackage{xcolor}
\usepackage[margin=0.85in]{geometry}
\usepackage{longtable,booktabs,array}
\usepackage{calc} % for calculating minipage widths
% Correct order of tables after \paragraph or \subparagraph
\usepackage{etoolbox}
\makeatletter
\patchcmd\longtable{\par}{\if@noskipsec\mbox{}\fi\par}{}{}
\makeatother
% Allow footnotes in longtable head/foot
\IfFileExists{footnotehyper.sty}{\usepackage{footnotehyper}}{\usepackage{footnote}}
\makesavenoteenv{longtable}
\setlength{\emergencystretch}{3em} % prevent overfull lines
\providecommand{\tightlist}{%
  \setlength{\itemsep}{0pt}\setlength{\parskip}{0pt}}
\setcounter{secnumdepth}{-\maxdimen} % remove section numbering
\usepackage{amsmath,amssymb}
\usepackage{microtype}
\setlength{\emergencystretch}{3em}
\ifLuaTeX
  \usepackage{selnolig}  % disable illegal ligatures
\fi
\usepackage{bookmark}
\IfFileExists{xurl.sty}{\usepackage{xurl}}{} % add URL line breaks if available
\urlstyle{same}
\hypersetup{
  pdftitle={Independently routed inverse sums and common rectangle obstructions},
  colorlinks=true,
  linkcolor={black},
  filecolor={Maroon},
  citecolor={Blue},
  urlcolor={blue},
  pdfcreator={LaTeX via pandoc}}

\title{Independently routed inverse sums and common rectangle
obstructions}
\usepackage{etoolbox}
\makeatletter
\providecommand{\subtitle}[1]{% add subtitle to \maketitle
  \apptocmd{\@title}{\par {\large #1 \par}}{}{}
}
\makeatother
\subtitle{Consolidated mathematical working draft}
\author{}
\date{11 October 2026}

\begin{document}
\maketitle

\section{Purpose and main
conclusions}\label{purpose-and-main-conclusions}

This manuscript gives a self-contained consolidation of an existing
family of additive representation obstructions. It establishes one
common low-rank rectangle for independently routed branches, then
compares that rectangle with a Fourier matrix or a normalized complex
flat unitary. There are two distinct mechanisms: componentwise
projective directions over an exact field, and weighted support
separators followed by a complex singular-value comparison. An earlier
ordered-separation mechanism is included to make the common-rectangle
and spectral statements precise and to explain which specialization the
separator theorem extends.

The principal complex theorem concerns \[
 M=\delta D+E+K+\sum_{i=1}^{r}L_iQ_i^{-1}R_i,\qquad d=s+1+\delta.
\] Here \(D\) is diagonal, \(E\) has at most \(s\) nonzeros in every
row, \(\operatorname{rank}K\le t\), and \(\delta\in\{0,1\}\). Each
\(Q_i\) is invertible, its symmetric bipartite support lift has marked
weighted half-separators of size at most \(q\), and each endpoint has
permutation-contained support, possibly with zero weights. Every branch
may have unrelated left and right routes. For a complex unitary \(U\)
with \(|U_{ab}|=N^{-1/2}\), exact equality \(U=M\) implies \[
 \boxed{N\le d^2(qr+t)9^r+dq(3^r-1)+2(d-1).}                 \tag{1}
\] The theorem requires integers \(q,r\ge1\) and \(s,t\ge0\). Invertible
column-two denominators permit \(q=2\). With fixed \(q,s,t\), (1) forces
\(r=\Omega(\log N)\) in any exact representation family.

A separately scoped algebraic theorem applies to the unnormalized
Fourier matrix \(F_{ab}=\omega^{ab}\) over a field containing a root
\(\omega\) of exact order \(N\). For invertible column-two denominators
and a row-degree-\(\delta_0\) extra summand \(S\), \[
 F=S+\sum_{i=1}^{r}L_iQ_i^{-1}R_i
 \quad\Longrightarrow\quad
 \boxed{N\le\delta_0(r+1)+r(r+1)^r.}                         \tag{2}
\] This theorem does not carry an uncharged low-rank correction, and it
uses neither positivity nor arbitrary-row full spark. It applies at
composite lengths and in positive characteristic whenever the required
root exists.

All results here are restricted additive representation statements. The
paper does not establish an unrestricted Fourier arithmetic lower bound,
a product-grammar obstruction, an optimal constant, a practical speedup,
or worldwide priority. The included results were already accepted within
the project's stated proof scopes. This consolidated presentation is a
working draft awaiting its own independent review; it is not a
publication-acceptance claim or a new result count. Exact historical
priority remains open.

\section{Definitions and quantifiers}\label{definitions-and-quantifiers}

\subsection{Coordinates and endpoints}\label{coordinates-and-endpoints}

Write \([N]=\{0,\ldots,N-1\}\), where \(N\ge1\). Submatrices use their
displayed row and column sets. Complements are taken in the full
coordinate set of the matrix in question. All ranks in the algebraic
statements are over the stipulated field.

A matrix has \textbf{permutation-contained support} if its support is
contained in a permutation matrix. Equivalently, it is a diagonal matrix
times a permutation matrix; the diagonal is allowed to be singular. For
a routed branch \(B_i=L_iQ_i^{-1}R_i\), choose permutations \(p_i,s_i\)
and endpoint weights so that \[
 B_i[a,b]=\ell_{i,a}\rho_{i,b}(Q_i^{-1})[p_i(a),s_i(b)].       \tag{3}
\] For definiteness, \(L_i[a,p_i(a)]=\ell_{i,a}\) and
\(R_i[s_i(b),b]=\rho_{i,b}\). This entrywise convention avoids a hidden
transpose convention. Zero entries of a partial permutation support can
be completed arbitrarily to a permutation and assigned zero weight.

The \(2r\) permutations in (3) are independent. They need not match
within a branch or across branches. Denominators, endpoint values, and
routes may all depend on \(N\). Coefficient magnitudes and condition
numbers are unrestricted. Invertibility is required of each denominator,
but not of its principal blocks or its endpoints.

\subsection{Exact fields and normalized
targets}\label{exact-fields-and-normalized-targets}

In the exact-field theorem, \(\mathbb F\) is any field and
\(\omega\in\mathbb F^\times\) has exact multiplicative order \(N\). Thus
the powers \(1,\omega,\ldots,\omega^{N-1}\) are distinct. If the
characteristic is positive, it cannot divide \(N\). The Fourier target
is \(F_{ab}=\omega^{ab}\) without a normalization requiring a square
root.

In the spectral theorem the field is \(\mathbb C\). A \textbf{normalized
flat unitary} means \(U^*U=I\) and \(|U_{ab}|=N^{-1/2}\) for every
entry. This includes normalized complex Fourier matrices at every length
and normalized complex Hadamard matrices. The symbol \(\|\cdot\|_2\)
denotes operator norm and \(\|\cdot\|_F\) Frobenius norm. Neither norm
statement is asserted over a finite field.

\subsection{The marked separator
hypothesis}\label{the-marked-separator-hypothesis}

For invertible \(Q\), define \[
 H_Q=\begin{pmatrix}0&Q\\Q^T&0\end{pmatrix},\qquad
 H_Q^{-1}=\begin{pmatrix}0&Q^{-T}\\Q^{-1}&0\end{pmatrix}.      \tag{4}
\] Its undirected support graph has an upper vertex for each row of
\(Q\) and a lower vertex for each column of \(Q\). The transpose is
ordinary transpose; \(H_Q\) need not be Hermitian.

Here is a sufficient hypothesis with explicit quantifiers. For every
equal-size pair of target sets \(A,B\subseteq[N]\) used to restrict the
branch, mark the lower vertices \(p_i(A)\) red and the upper vertices
\(s_i(B)\) blue, with unit vertex weight on every mark and zero
elsewhere. There must exist a vertex set \(X_i\) of size at most \(q\)
whose deletion leaves every connected component with at most half the
total marked weight. Red and blue marks lie in different halves even if
their numeric coordinate labels agree. Mark endpoint positions even when
the associated endpoint weight is zero.

It suffices to have this property for all such equal-size pairs; more
minimally, it suffices at every marking actually encountered by the
construction, with a valid choice at every stage. We use the all-pairs
sufficient formulation in theorem statements so that the quantifier is
independent of choices made during a proof. The stronger, familiar
condition requiring a half-component separator for \textbf{every
nonnegative vertex weighting} also suffices. Treewidth at most \(q-1\)
is sufficient, but is not the definition of the hypothesis.

Separators, component groupings, and orientations may be chosen afresh
for each branch and each current rectangle. No common graph, common tree
decomposition, or common coordinate ordering is assumed.

\subsection{Correction conventions}\label{correction-conventions}

A row-degree bound is a maximum over individual rows. A bound on average
degree or total nonzeros is not the same hypothesis. In (1), \(E\) has
row degree at most \(s\); its column degrees and values are
unrestricted. The optional diagonal is charged separately by \(\delta\).
The rank-\(t\) correction \(K\) is not removed by sparsity deletion; it
survives and is charged to the final rank budget.

In (2), \(S\) has row degree at most \(\delta_0\). A free diagonal can
be included by taking \(\delta_0=1\). This field theorem has no separate
\(K\) allowance. The parameters \(\delta_0\), \(\delta\), and \(d\) must
not be identified.

\section{Common rectangles from ordered
separation}\label{common-rectangles-from-ordered-separation}

This section records the existing ordered specialization. Its algebraic
rectangle selection and its complex comparison have different field
requirements. It is useful when the branches come with compatible
low-rank separated rectangles in their own orders; the later graph
theorem does not require such a global order.

\subsection{Theorem O for independently ordered
sums}\label{theorem-o-for-independently-ordered-sums}

Let \(b\ge1\) and let each \(\rho_i\) be a nonnegative integer. For
\(i=1,\ldots,b\), let \(B_i\) be a matrix over a field and let
\(a_i,c_i:[N]\to[N]\) be bijective row and column orderings. Suppose \[
 \operatorname{rank}B_i[A,B]\le \rho_i
\] whenever the nonempty sets \(A,B\) are strictly separated in either
orientation: \[
 \max a_i(A)<\min c_i(B),\quad\hbox{or}\quad
 \max c_i(B)<\min a_i(A).
\] Let \(h\) be the number of distinct specified order pairs, and put \[
 R_O=\sum_i\rho_i,\qquad m=\lfloor N/2^h\rfloor.
\] There exist target row and column sets of size \(m\) on which
\(\sum_iB_i\) has rank at most \(R_O\). Empty sets are permitted if
\(m=0\), making the conclusion vacuous.

\textbf{Proof.} For current row and column sets of equal size \(n\ge2\),
put \(u=\lfloor n/2\rfloor\). Let the low and high \(u\) rows be chosen
in the current row order, and similarly for columns. If low rows precede
high columns, keep them. Otherwise, if low columns also failed to
precede high rows, the inequalities would give \[
 \max a(A_{\rm lo})\ge\min c(B_{\rm hi})
 >\max c(B_{\rm lo})\ge\min a(A_{\rm hi})
 >\max a(A_{\rm lo}),
\] an impossibility. The opposite orientation therefore works. Apply
this once per distinct order pair. Strict separation already obtained
survives taking subsets. Repeated integer halving gives
\(m=\lfloor N/2^h\rfloor\). Rank subadditivity proves the claim. If a
stage reaches size zero or one before all pairs have been processed, the
final guaranteed size is zero and the empty rectangle suffices.
\(\square\)

\subsection{Weighted cycle Green branches satisfy the
hypothesis}\label{weighted-cycle-green-branches-satisfy-the-hypothesis}

Let \(W[v_j,v_{j+1}]=w_j\) on a specified directed Hamiltonian cycle,
with indices modulo \(N\), and let all other entries be zero. Put
\(\lambda=\prod_jw_j\). Then \(W^N=\lambda I\), and for \(\lambda\ne1\),
\[
 (I-W)^{-1}=\frac{I+W+\cdots+W^{N-1}}{1-\lambda}.
\] An entry is the product on the unique directed arc of length
\(0,\ldots,N-1\), divided by \(1-\lambda\). Between strictly separated
row and column positions, this product factors at a cut into a row
factor and a column factor. In the reverse orientation it factors across
the cycle wrap. No individual edge weight is divided out. Thus zero
edges are allowed and each separated rectangle has rank at most one.
Independent endpoint routes transport the two orders separately;
endpoint scales, including zero scales, cannot increase rank.

A direct sum of legal cycle blocks has the same property when the block
orders are concatenated: at most one component can contribute a nonzero
interaction crossing a separating cut. Arbitrarily many bare cycles
therefore do not by themselves require the later pseudoforest extension.
Attached trees and unbounded row degree are genuinely outside this
immediate cycle-order argument.

\subsection{Corollary O for a normalized complex
target}\label{corollary-o-for-a-normalized-complex-target}

For a normalized flat unitary and \(T=\sum_iB_i\), Theorem O gives \[
 U=T\ \Longrightarrow\ m^2\le NR_O,
\] and, for \(m>0\), \[
 \|U-T\|_F^2\ge\max(0,m^2/N-R_O),\qquad
 \|U-T\|_2^2\ge\max(0,m/N-R_O/m).                           \tag{5}
\] For \(m=0\) the lower bounds are read as zero. The rank and norm
proof is given in the next section. With \(b\) rank-one branches,
\(h\le b\); fixed \(b\) gives asymptotic operator-error floor at least
\(2^{-b/2}\). The associated logarithmic exact branch obstruction was
already present in this ordered setting. The \(9^r\) theorem below
should not be presented as a uniform numerical improvement over this
\(2^h\) retention theorem on Hamiltonian specializations. We do not need
the separate prime-complex full-spark refinement here.

\section{Rank and spectral comparison
lemmas}\label{rank-and-spectral-comparison-lemmas}

\subsection{Lemma N for complementary inverse
blocks}\label{lemma-n-for-complementary-inverse-blocks}

For an invertible \(n\times n\) matrix \(A\) over any field and
arbitrary index sets \(I,J\), \[
 \operatorname{rank}A^{-1}[I,J]
 =\operatorname{rank}A[J^c,I^c]+|I|+|J|-n.                 \tag{6}
\] \textbf{Proof.} Take \(y\) supported on \(J\) and set \(z=A^{-1}y\).
The equation \(z|_I=0\) is equivalent to \(z\) being supported on
\(I^c\) and \((Az)|_{J^c}=0\). The map \(y\mapsto z\) bijects the
corresponding kernels. Their domain dimensions are \(|J|\) and
\(n-|I|\); rank-nullity gives (6). The complements are swapped. No
principal-block inverse is used. \(\square\)

If a vertex separator \(X\) partitions the remaining support graph of an
invertible symmetric-support matrix \(A\) into component unions
\(V_1,V_2\), then \(A[V_1,V_2]=0\). The complementary block \[
 A[V_2^c,V_1^c]=A[V_1\cup X,V_2\cup X]
\] has rank at most \(2|X|\): its entries outside the zero
\(V_1\times V_2\) block occupy \(|X|\) extra rows or \(|X|\) extra
columns. Since \(|V_1|+|V_2|=n-|X|\), (6) gives \[
 \operatorname{rank}A^{-1}[V_1,V_2]\le |X|.                \tag{7}
\] The same is true with the sides exchanged. Applying this to (4) and
restricting to lower rows and upper columns proves the inverse-rectangle
rank bound for \(Q^{-1}\). Endpoint scalings preserve the upper bound
even when they vanish.

\subsection{Lemma S for a low-rank rectangle of a flat
unitary}\label{lemma-s-for-a-low-rank-rectangle-of-a-flat-unitary}

Let \(V=U[A,B]\) with \(|A|=|B|=m\). Then \[
 \|V\|_F^2=m^2/N,\qquad \|V\|_2\le1,\qquad
 \operatorname{rank}V\ge m^2/N.                            \tag{8}
\] The first identity uses flat entries and the second uses coordinate
compression of a unitary. The last follows by summing squared singular
values, each at most one.

If \(W=M[A,B]\) has rank at most a nonnegative integer \(R\), then \[
 \|U-M\|_F^2\ge\max(0,m^2/N-R).                            \tag{9}
\] For \(m>R\), \[
 \|U-M\|_2^2\ge
 \max\left(0,\frac{m^2/N-R}{m-R}\right).                  \tag{10}
\] For completeness, the squared distance from \(V\) to any
rank-at-most-\(R\) matrix is at least the sum of the squared singular
values after the first \(R\). One can see this by projecting onto the
orthogonal complement of the candidate's column space: a rank-\(R\)
orthogonal projector captures at most the sum of the top \(R\) squared
singular values. Those top values contribute at most \(R\) by (8),
proving (9). If the block's operator error is \(\varepsilon\), its
\((R+1)\)st singular value is at most \(\varepsilon\): the span of its
top \(R+1\) right singular vectors intersects the kernel of \(W\)
nontrivially. Every later singular value is then also at most
\(\varepsilon\). There are \(m-R\) such values, giving \[
 m^2/N\le R+(m-R)\varepsilon^2.
\] Coordinate restriction does not increase either norm, proving (10).
The weaker (5) follows from (9) and \(\|V-W\|_F^2\le m\|V-W\|_2^2\). If
\(R\ge m\), the Frobenius bound is zero and (10) is not used.
\(\square\)

\section{Separator selection and the complex corrected
theorem}\label{separator-selection-and-the-complex-corrected-theorem}

\subsection{Lemma C for two colors}\label{lemma-c-for-two-colors}

Suppose a branch satisfies the marked separator hypothesis, and the
current target rectangle has \(m\) rows and \(m\) columns with
\(m\ge4q\). There are subsets of equal size \[
 m'=\left\lceil\frac{m-q}{3}\right\rceil                  \tag{11}
\] on which that branch has rank at most \(q\).

\textbf{Proof.} Mark the routed positions as in the definition and
remove a half-separator \(X\), \(|X|\le q\). Each color has at least
\(m-q\) surviving marks. Trim each color to exactly \(u=m-q\).
Components have at most \(m\) retained marks, and \[
 m\le4u/3,\qquad\hbox{total retained marks}=2u.
\] Group components into two sides with each side's retained weight
between \(2u/3\) and \(4u/3\). If a component has weight at least
\(2u/3\), take that component for the first side; it has weight at most
\(4u/3\). Otherwise accumulate components until that side first reaches
\(2u/3\); it remains below \(4u/3\). The complement is valid in either
case. Place unmarked components arbitrarily so the two sides cover all
vertices outside \(X\).

If the first side has \(a\) red and \(b\) blue marks, the second has
\(u-a\) red and \(u-b\) blue marks. At least one oriented cross pair \[
 (a,u-b),\qquad (u-a,b)
\] has both counts at least \(u/3\). If both failed, combining their
failure disjunctions would give either a side of total weight below
\(2u/3\), or \(a<u/3\) and \(u-a<u/3\), or \(b<u/3\) and \(u-b<u/3\).
Each is impossible. Integer counts therefore permit \(\lceil u/3\rceil\)
rows and columns. They lie on opposite component unions, so (7) and
endpoint rank monotonicity give the claimed rank at most \(q\).
\(\square\)

The retained fraction \(1/3\) is sharp for this component-count
information: when \(u\) is divisible by three, three components each
with \(u/3\) marks of each color allow no larger oriented cross minimum.
This observation does not prove the final \(9^r\) exponent or constants
optimal for matrices.

\subsection{Lemma R for removing sparse
corrections}\label{lemma-r-for-removing-sparse-corrections}

Put \(d=s+1+\delta\) and \(p=\lfloor N/d\rfloor\). Choose any \(p\)
target rows \(A_0\). The union of column supports of \(E\) in these
rows, together with \(A_0\) itself if \(\delta=1\), has size at most \[
 (s+\delta)p=(d-1)p.
\] Since \(N-(d-1)p\ge p\), choose \(p\) columns \(B_0\) outside that
union. Both \(E[A_0,B_0]\) and \(\delta D[A_0,B_0]\) vanish. This uses
only row degree. The sets need not be disjoint when \(\delta=0\). The
low-rank term satisfies \(\operatorname{rank}K[A_0,B_0]\le t\) and is
retained. These facts persist under further restrictions. \(\square\)

\subsection{Theorem C for correction compatible inverse
sums}\label{theorem-c-for-correction-compatible-inverse-sums}

Assume the complex model and all hypotheses in the opening statement,
including \(q,r\ge1\) and the all-pairs marked separator condition. Then
(1) holds.

\textbf{Proof.} Begin with the rectangle in Lemma R and process the
\(r\) branches in any fixed order. Let \[
 m_0=p,\qquad m_{j+1}=\left\lceil\frac{m_j-q}{3}\right\rceil.\tag{12}
\] Whenever each preceding \(m_j\ge4q\), Lemma C constructs nested
target rectangles with these sizes. A rank bound established for an
earlier branch persists under every subsequent restriction. On the final
square, each inverse branch has rank at most \(q\), sparse and diagonal
corrections vanish, and \(K\) has rank at most \(t\). Thus the rank
budget is \[
 R=qr+t.                                                   \tag{13}
\] Ceiling monotonicity yields, for every legal stage, \[
 m_j\ge\frac p{3^j}-\frac q2(1-3^{-j}).                     \tag{14}
\] Put \[
 a=3^r,\qquad c=(d-1)+\frac{dq(a-1)}2,\qquad B=d^2Ra^2.
\] Assume for contradiction that \(N>B+2c\). We verify legality before
relying on the final recurrence. At the initial stage and, inductively,
at each subsequent stage \(j\le r-1\), (14) and \(p\ge(N-d+1)/d\) give
\[
 \begin{split}
 m_j&\ge\frac{N-d+1}{d3^j}-\frac q2(1-3^{-j})\\
    &>dR3^{2r-j}-1-q/2\\
    &\ge9q-1-q/2\ge4q.
 \end{split}                                               \tag{15}
\] Here \(N>B\), \(d\ge1\), \(R\ge qr\ge q\), \(r\ge1\), and
\(2r-j\ge r+1\ge2\). The strict threshold therefore guarantees every
required use of Lemma C without circularity.

At the final stage, retaining the positive term in (14) gives \[
 m=m_r\ge\frac{N-c}{da}>0.                                 \tag{16}
\] If \(U=M\), (8) and (13) imply \(m^2\le NR\), and hence \[
 (N-c)^2\le NB.
\] But \[
 (N-c)^2-NB=N\bigl[N-(B+2c)\bigr]+c^2>0,
\] a contradiction. Expanding \(B+2c\) gives (1). \(\square\)

\subsection{Base specializations and quantitative
form}\label{base-specializations-and-quantitative-form}

The uncorrected case \(s=t=\delta=0\) gives \[
 N\le qr9^r+q(3^r-1).                                      \tag{17}
\] With only an arbitrary diagonal, \(s=t=0\), \(\delta=1\), \[
 N\le4qr9^r+2q(3^r-1)+2.                                   \tag{18}
\] For example, \(q=2,r=2,s=1,t=1,\delta=1\) gives \(N\le3697\).

For approximation, exact equality is unnecessary. For any \(N\) at which
(12) has \(m_j\ge4q\) before every selection, the construction produces
a block with rank at most \(R=qr+t\). Equations (9) and (10) therefore
apply to \(M\), with \(m=m_r\). A sufficient condition for legality is
violation of the exact cutoff, as proved in (15); legality may also be
checked directly at smaller sizes. If a step is illegal, this
construction by itself supplies no claimed nonempty rectangle or
positive error floor.

With \(q,r,s,t\) fixed and \(N\to\infty\), (12) gives \[
 \frac{m_r}{N}\longrightarrow \alpha=\frac1{d3^r}.
\] Consequently, uniformly over the admissible choices of matrices and
routes at each size, \[
 \liminf\|U-M\|_2\ge\sqrt\alpha,\qquad
 \liminf\frac{\|U-M\|_F}{\|U\|_F}\ge\alpha.                \tag{19}
\] Here \(\|U\|_F=\sqrt N\). Fixed branch count cannot yield vanishing
error in these normalized metrics. The floors decay as \(r\) grows; they
do not imply \(r=\Omega(\log N)\) for an arbitrary fixed positive error
tolerance.

For exact equality with fixed \(q,s,t\), the right side of (1) is at
most \(C r9^r\) for a constant \(C\) depending only on those fixed
parameters. Taking logarithms and using \(\log r\le r\) for \(r\ge1\)
proves the stated \(\Omega(\log N)\) branch lower bound. This asymptotic
statement does not apply unchanged when the correction or separator
budgets grow with \(N\).

\section{Sufficient support conditions including column
two}\label{sufficient-support-conditions-including-column-two}

\subsection{Tree decompositions}\label{tree-decompositions}

A tree decomposition has vertex bags that cover every support edge and
whose occurrences of each graph vertex form a connected subtree. If
every bag has at most \(q\) vertices, assign each graph vertex's
nonnegative weight to one bag containing it. A finite weighted tree has
a centroid: repeatedly move toward a component of weight greater than
half the total, if one exists. No traversed edge can be crossed back in
that direction, so this ends. Remove the centroid bag. Every remaining
graph component lies entirely on one side of the bag tree and
consequently has at most half the weight. This proves the sufficient
separator property, also for disconnected graphs by joining component
decomposition trees without adding bag vertices. Zero total weight is
trivial.

\subsection{Lemma P for an invertible column two
denominator}\label{lemma-p-for-an-invertible-column-two-denominator}

Suppose every column of an invertible \(N\times N\) matrix \(Q\) has at
most two actual nonzero entries. A zero column is impossible. Form its
row-support multigraph: row coordinates are vertices, two-entry columns
are labeled edges, and singleton columns are labeled loops. For a
connected component with \(v\) vertices and \(e\) column-edges, column
independence gives \(e\le v\). Both total counts are \(N\), so equality
holds componentwise. Each component is therefore unicyclic, allowing a
loop or two parallel edges as its cycle.

In the bipartite lift, a two-entry column subdivides its edge; a
singleton column replaces its loop by a leaf. Thus every component of
the lifted graph is a tree or unicyclic. It is a pseudoforest. Unbounded
row degree, many components, and singleton columns are allowed.

\subsection{Weighted pseudoforests permit two separator
vertices}\label{weighted-pseudoforests-permit-two-separator-vertices}

Give a pseudoforest arbitrary nonnegative weights of total \(W\). If
\(W=0\), or all components have weight at most \(W/2\), the empty
separator works. Otherwise there is a unique heavy component. A heavy
tree has a one-vertex centroid with component weights at most half that
tree's weight, hence at most \(W/2\).

For a heavy unicyclic component, let \(T_v\) be the tree attached to
cycle vertex \(v\), including \(v\), and let its weight be \(w_v\). If
some \(w_v>W/2\), add the outside weight \(W-w_v\) at the root \(v\) in
\(T_v\) and choose a centroid \(z\) of that augmented tree. If
\(z\ne v\), the part containing \(v\) can reconnect to the outside
cycle, but its actual weight is at most the augmented part's weight, at
most \(W/2\). Every other detached part has the same weight as in the
augmented tree. If \(z=v\), the cycle remainder weighs at most
\(W-w_v<W/2\), and the detached branches satisfy the centroid bound.
Other original graph components were already light.

If instead all \(w_v\le W/2\), delete any cycle vertex \(x\). Aggregate
each remaining attached tree into its cycle root, producing a weighted
path. Choose a centroid \(y\) of that path. Its two path sides have
aggregate weight at most half the remaining cycle-component weight,
hence at most \(W/2\). Branches detached at \(x\) or \(y\) have weight
at most \(w_x\) or \(w_y\), hence at most \(W/2\). Thus \(\{x,y\}\)
suffices. If a remaining path has zero weight, any of its vertices is a
valid centroid. A genuine cycle in the bipartite lift has at least four
vertices; singleton-column loops have already become tree components.

Therefore Theorem C permits \(q=2\) for every invertible column-two
denominator. A forest lift permits \(q=1\). Merely assuming a fixed
column degree greater than two gives no bounded-separator guarantee.
Weighted-permutation Green denominators \(I-W\), when invertible, have
at most two actual nonzeros per column and are included, with zero edge
weights allowed. \(\square\)

\section{The separate exact field
theorem}\label{the-separate-exact-field-theorem}

\subsection{Lemma F for inverse row
directions}\label{lemma-f-for-inverse-row-directions}

Let \(Q\) be invertible and column-two over a field, and select any set
\(J\) of \(k\) rows of \(Q^{-1}\), \(1\le k\le N\). The nonzero columns
of \(A=Q^{-1}[J,:]\) occupy exactly \(k\) distinct projective directions
with linearly independent representatives.

\textbf{Proof.} Delete the \(k\) column-labeled edges in \(J\) from the
row-support multigraph of Lemma P. This is deletion of denominator
\textbf{columns}, not denominator rows. The resulting graph is a
pseudoforest, with exactly \(k\) tree components because \[
 \sum_C\bigl(|V(C)|-|E(C)|\bigr)=N-(N-k)=k,
\] and each tree contributes one while each unicyclic component
contributes zero. Isolated vertices count as trees. A surviving
unicyclic component need not be untouched: deleting an attached-tree
edge can split off a tree while retaining a smaller unicyclic component.

Put \(B=Q[:,J^c]\). Its \(N-k\) columns remain independent, also when
restricted to their support components. A tree with \(v\) vertices has
\(v-1\) such columns and local left nullity one; a unicyclic component
has \(v\) columns and local left nullity zero. Hence \[
 W=\{x\in\mathbb F^{1\times N}:xB=0\}
\] is a direct sum of \(k\) one-dimensional spaces with disjoint vertex
supports. The rows of \(A\) are independent, lie in \(W\) by
\(Q^{-1}Q=I\), and form a basis of \(W\).

Choose one nonzero local null vector for each tree and stack them into a
matrix \(H\). There is an invertible \(k\times k\) matrix \(C\) with
\(A=CH\). A column corresponding to a vertex in tree \(u\) is a scalar
multiple of \(C[:,u]\); vertices in the unicyclic components give zero
columns. Every local null vector has a nonzero coordinate, so every one
of the \(k\) independent directions appears. No division by an edge
weight and no characteristic restriction has been used. \(\square\)

\subsection{Routing and common
refinement}\label{routing-and-common-refinement}

Restrict a routed branch to any \(k\) target rows. Choose the \(k\)
distinct denominator inverse rows prescribed by its left permutation,
including positions with zero endpoint weight. Lemma F supplies \(k\)
directions; multiplying by the diagonal endpoint may collapse some to
zero but creates none. The right endpoint permutes and scales columns.
Therefore the restricted target columns admit a partition into at most
\(k\) classes, each of rank at most one. Absorb zero columns into any
class; an all-zero branch uses one class. A zero vector is not a
projective point, but does not need a separate linear-rank class.

The partitions from different branches are unrelated. Their common
refinement has at most \(k^r\) cells, and on each cell the sum of the
\(r\) branches has rank at most \(r\). This is a family of common
rectangles on a fixed set of rows, distinct from the one nested square
selected by graph separators.

\subsection{Theorem F with a row sparse
summand}\label{theorem-f-with-a-row-sparse-summand}

Let \(F_{ab}=\omega^{ab}\) over a field with \(\omega\) of exact order
\(N\). Suppose \[
 F=S+\sum_{i=1}^{r}L_iQ_i^{-1}R_i,
\] where \(r\ge1\), every \(Q_i\) is invertible and column-two,
endpoints satisfy (3), and \(S\) has row degree at most
\(\delta_0\ge0\). Then (2) holds.

\textbf{Proof.} If \(r\ge N\), the claimed inequality is automatic.
Otherwise put \(k=r+1\le N\) and select consecutive Fourier rows
\(0,\ldots,k-1\). At most \(\delta_0k\) columns have a nonzero \(S\)
entry on these rows; discard them. Refine the \(r\) branch partitions on
the remaining columns as above.

Any \(k\) distinct Fourier columns on the selected rows form a
Vandermonde matrix with determinant \[
 \prod_{u<v}(\omega^{b_v}-\omega^{b_u})\ne0.
\] Thus a common-refinement cell cannot contain \(k\) surviving columns,
since the branch sum there has rank at most \(r=k-1\). Every cell has at
most \(r\) columns, so at most \(rk^r\) columns survive. Restoring the
discarded columns proves \(N\le\delta_0k+rk^r\). \(\square\)

Any consecutive row interval, including a cyclic wrap, works by nonzero
column scaling. The proof does not allow arbitrary selected rows at
composite length. It gives \(N\le r(r+1)^r\) without \(S\), and
\(N\le(r+1)+r(r+1)^r\) with a free diagonal. In particular two
uncorrected branches are excluded for \(N\ge19\); no attainability claim
is made below this cutoff. The uncorrected field bound yields
\(r=\Omega(\log N/\log\log N)\), weaker than the complex separator
asymptotic but with a wider field scope. For fixed \(\delta_0\), the
same field asymptotic follows.

If \(r=0\), the theorem's proof is not invoked: uncorrected equality to
the zero matrix is impossible for \(N\ge1\); with \(S\) alone, all
Fourier entries are nonzero, so equality requires \(\delta_0\ge N\). A
diagonal-only exact representation is possible at \(N=1\) and impossible
at \(N>1\). For the complex model with other sparse or rank corrections
and \(r=0\), Theorem C makes no claim; those corrections may themselves
represent the target.

\section{Grammar and quantifier
concordance}\label{grammar-and-quantifier-concordance}

\begin{longtable}[]{@{}
  >{\raggedright\arraybackslash}p{(\columnwidth - 6\tabcolsep) * \real{0.2500}}
  >{\raggedright\arraybackslash}p{(\columnwidth - 6\tabcolsep) * \real{0.2500}}
  >{\raggedright\arraybackslash}p{(\columnwidth - 6\tabcolsep) * \real{0.2500}}
  >{\raggedright\arraybackslash}p{(\columnwidth - 6\tabcolsep) * \real{0.2500}}@{}}
\toprule\noalign{}
\begin{minipage}[b]{\linewidth}\raggedright
Result
\end{minipage} & \begin{minipage}[b]{\linewidth}\raggedright
Algebra and target
\end{minipage} & \begin{minipage}[b]{\linewidth}\raggedright
Branch and correction contract
\end{minipage} & \begin{minipage}[b]{\linewidth}\raggedright
Guaranteed conclusion
\end{minipage} \\
\midrule\noalign{}
\endhead
\bottomrule\noalign{}
\endlastfoot
Theorem O & Any field; arbitrary matrix sum & Specified independent row
and column order pairs; separated rank budgets \(\rho_i\) & One square
of side \(\lfloor N/2^h\rfloor\); rank at most \(\sum_i\rho_i\) \\
Corollary O & Complex normalized flat unitary & Same ordered sum; no
imported sparse-correction rule & Exact rank comparison and (5) \\
Theorem C & Complex normalized flat unitary & Invertible denominators
with marked half-separators; independent permutation-contained
endpoints; \(\delta D+E+K\) & Cutoff (1), with the legal-recurrence norm
bounds (9)--(10) \\
Column-two specialization & Complex normalized flat unitary & Any
invertible actual column-two denominators, including attached trees &
Theorem C with \(q=2\); \(q=1\) for forest lifts \\
Theorem F & Exact-order-root Fourier matrix over a field & Invertible
column-two denominators; independent permutation-contained endpoints;
row-degree-\(\delta_0\) term only & Bound (2); no singular-value or
uncharged rank-correction statement \\
\end{longtable}

In Theorems C and F the universal quantifier is over all admissible
matrices, routes, weights, and corrections at the specified size. The
existential objects are the retained rectangle, separators or
direction-class partitions. They may depend on the entire proposed
representation. No fixed rectangle is asserted to work for all
representations simultaneously.

Theorem O counts distinct specified order pairs \(h\) separately from
the number of branches. Theorem C processes \(r\) inverse branches and
does not silently replace \(r\) by a number of graph templates. Theorem
F counts common partition cells, not graph separators. All three
constructions use rank monotonicity under restriction, but their
retention formulas and hypotheses are not interchangeable.

A sparse matrix is eliminated by a support argument, whereas a low-rank
matrix is charged by rank subadditivity. One cannot substitute one
correction convention for the other or transfer the \(t\) term in
Theorem C into (2) without a separately proved statement. No such
extension is asserted here.

\section{Scaling and edge cases}\label{scaling-and-edge-cases}

Independent \textbf{target} row or column permutations preserve flat
unitarity, so Theorem C applies directly to such normalized targets.
Unitary diagonal phases also preserve it. Arbitrary nonunitary diagonal
target scaling does not preserve the numerical normalization in (8).

The accepted scaling rule can nevertheless be stated explicitly. If the
target is \(A=D_LUD_R\) with invertible diagonal \(D_L,D_R\), transform
an approximant to \(\widetilde M=D_L^{-1}MD_R^{-1}\). The endpoint
grammar, row degree, rank correction, and diagonality of the optional
diagonal are preserved. Exact branch-count obstructions therefore
transfer. For either norm, \[
 \|A-M\|\ge d_{\min}e_{\min}\,\|U-\widetilde M\|,
\] where \(d_{\min}=\min|D_L[a,a]|\) and \(e_{\min}=\min|D_R[b,b]|\).
This follows by applying the norm inequality to multiplication by the
inverse diagonals. Thus an absolute error floor must be multiplied by
these minima. Relative operator floors can be weakened by the factor
\(\kappa_2(D_L)\kappa_2(D_R)\). A condition number alone does not set an
absolute scale; even scalar rescaling changes it. Singular target
scaling and arbitrary dense preconditioners are outside this transfer
rule.

Other safeguards are worth collecting:

\begin{itemize}
\tightlist
\item
  Empty separators, zero graph weights, zero endpoints, and disconnected
  supports are permitted. The proof marks zero endpoint positions
  conservatively.
\item
  Integer ceilings and floors are part of the statements. No asymptotic
  argument is used to justify a finite recurrence step.
\item
  The assumption \(q\ge1\) avoids a separate zero-separator convention.
  The results do not assert a zero right-hand side for a zero-branch
  diagonal case.
\item
  Column-two refers to actual nonzero support. Vanishing coefficients
  are handled before the support graph is formed. In a formal
  permutation-superpattern proof one must instead keep formal edges and
  redo the appropriate count; these are different bookkeeping methods.
\item
  Equality at a necessary cutoff is not construction evidence. The
  constants for \(r\ge2\) and the separator exponent are not claimed
  optimal.
\item
  A required branch count and an available branch-by-branch evaluation
  schedule do not prove a lower bound on every shared evaluator.
  Products, interleaved sums, nonmonomial endpoint mixing, and larger
  additive grammars remain excluded.
\end{itemize}

\section{Relation to established
literature}\label{relation-to-established-literature}

The following comparison uses the primary passages already identified by
the bounded source audit. The references were reopened for this
consolidation; this was a citation check, not a new exhaustive priority
search. The proofs above are written out so the mathematical conclusions
do not depend on accepting a negative literature-search claim.

\begin{longtable}[]{@{}
  >{\raggedright\arraybackslash}p{(\columnwidth - 4\tabcolsep) * \real{0.3333}}
  >{\raggedright\arraybackslash}p{(\columnwidth - 4\tabcolsep) * \real{0.3333}}
  >{\raggedright\arraybackslash}p{(\columnwidth - 4\tabcolsep) * \real{0.3333}}@{}}
\toprule\noalign{}
\begin{minipage}[b]{\linewidth}\raggedright
Ingredient or quantifier
\end{minipage} & \begin{minipage}[b]{\linewidth}\raggedright
Established primary result
\end{minipage} & \begin{minipage}[b]{\linewidth}\raggedright
Relation to this manuscript
\end{minipage} \\
\midrule\noalign{}
\endhead
\bottomrule\noalign{}
\endlastfoot
Swapped complementary inverse ranks & Vandebril and Van Barel, Section
3, Corollary 1 & Exact classical identity behind (6); the separator rank
bound is a short consequence \\
Rank structure induced by a graph & Chandrasekaran, Epperly, and
Govindarajan, Propositions 2.1--2.3 & GIRS gives inverse and sum/product
rank algebra on one graph; the present selection must handle unrelated
branch graphs \\
Tree rank profiles and algebra & Govindarajan, Chandrasekaran, and
Dewilde, Theorem 4.4, Corollary 4.5, Propositions 4.6--4.8 & TQS already
treats tree representations and compatible same-tree algebra, including
separate row and column block permutations \\
Weighted separator existence & Bodlaender et al., Section 2.1, Lemmas
2.1--2.2 and preceding weighted statement & Half-component separators
and scalar component grouping are classical; the red/blue iteration is a
specialized assembly \\
Componentwise projective directions & Zaslavsky, Section 2.1,
contraction map and Theorem 2.1(b) & Classical frame/gain-graph
antecedent of the inverse-row direction mechanism \\
The column-two class & Nelson and van Zwam, Section 2, frame-matrix
definition and Propositions 2.1--2.2 & Established matrix class; this
manuscript proves its needed rank-nullspace statement directly,
including characteristic two \\
\end{longtable}

For GIRS, the cited rank bound is expressed using border edges. A small
vertex separator is not interchangeable with a small number of border
edges when degrees are unbounded; (6)--(7) provide the needed
vertex-based argument. Common-graph addition alone also does not select
one rectangle across incompatible branch graphs. These are differences
between the inspected statements and the present contract, not a proof
that no derivation from GIRS is possible.

For TQS, it would be incorrect to claim that the theory allows only one
simultaneous row/column permutation. The final paper permits distinct
row and column block permutations. The present comparison concerns the
separate branch-specific routed partitions, common-rectangle selection,
and normalized target obstruction. No TQS representation problem,
optimal storage theorem, or published open question is claimed solved.

The frame/gain-graph literature already supplies the structural
perspective behind column-two matrices and componentwise directions. In
particular, the inverse-row argument is not the invention of a new
matrix class. The direct proof avoids importing a gain convention that
would mishandle singleton columns over characteristic two.
Consecutive-row Vandermonde nonvanishing and the singular-value energy
inequalities are also standard ingredients.

The bounded audit did not locate a theorem statement with the entire
independently routed, branch-specific marked-separator,
sparse-plus-low-rank-correction contract and exact cutoff (1). The
defensible description is a specialized synthesis/application with an
explicit quantified statement. That absence of a located match is not
proof of originality. Exact historical priority remains unresolved. The
earlier report's temporary absence of the correction author body is now
closed by the recovered matching body; this provenance correction does
not settle priority.

\subsection{Primary references}\label{primary-references}

\begin{enumerate}
\def\labelenumi{\arabic{enumi}.}
\tightlist
\item
  R. Vandebril and M. Van Barel, \emph{A short note on the nullity
  theorem}, KU Leuven report TW397, 15 July 2004. Section 3, Theorem 1,
  printed p.~4, and Corollary 1, printed p.~5.
  \href{https://www.cs.kuleuven.be/publicaties/rapporten/tw/TW397.pdf}{Institutional
  text}. This later primary account credits Fiedler--Markham and
  Gustafson; their original articles were not independently recovered in
  the prior bounded audit.
\item
  S. Chandrasekaran, E. N. Epperly, and N. Govindarajan,
  \emph{Graph-Induced Rank Structures and their Representations},
  arXiv:1911.05858v3. Propositions 2.1--2.3, printed pp.~5--8;
  high-degree Example 2.4.
  \href{https://arxiv.org/pdf/1911.05858v3}{Primary preprint}.
\item
  N. Govindarajan, S. Chandrasekaran, and P. Dewilde, \emph{Tree
  Quasi-Separable Matrices: A Simultaneous Generalization of
  Sequentially and Hierarchically Semiseparable Representations}, SIAM
  Journal on Matrix Analysis and Applications 46(2), 1562--1586, 2025,
  DOI 10.1137/24M1682348. Theorem 4.4 and following algebraic results,
  pp.~1575--1576.
  \href{https://repository.tudelft.nl/file/File_6b7549c8-222e-4316-8ac4-a113d7d4da1e}{Final
  institutional copy}.
\item
  H. L. Bodlaender, P. G. Drange, M. S. Dregi, F. V. Fomin, D.
  Lokshtanov, and M. Pilipczuk, \emph{A \(c^k n\) 5-Approximation
  Algorithm for Treewidth}, author version dated 5 February 2016,
  Section 2.1, pp.~5--6.
  \href{https://sites.cs.ucsb.edu/~daniello/papers/treewidthApproxJ.pdf}{Author
  text}.
\item
  T. Zaslavsky, \emph{Biased Graphs IV: Geometrical Realizations},
  Journal of Combinatorial Theory Series B 89 (2003), 231--297. Section
  2.1 in the author version, PDF pp.~8--10, especially the contraction
  map and Theorem 2.1(b).
  \href{https://people.math.binghamton.edu/zaslav/Tpapers/bg4.pdf}{Author
  text}. Author-PDF page numbers are not journal page numbers.
\item
  P. Nelson and S. H. M. van Zwam, \emph{On the existence of
  asymptotically good linear codes in minor-closed classes}, Section 2,
  pp.~4--5.
  \href{https://www.math.uwaterloo.ca/~apnelson/asympgood.pdf}{Author
  text}.
\end{enumerate}

\section{Theorem dependencies and proof source
lineage}\label{theorem-dependencies-and-proof-source-lineage}

\begin{longtable}[]{@{}
  >{\raggedright\arraybackslash}p{(\columnwidth - 4\tabcolsep) * \real{0.3333}}
  >{\raggedright\arraybackslash}p{(\columnwidth - 4\tabcolsep) * \real{0.3333}}
  >{\raggedright\arraybackslash}p{(\columnwidth - 4\tabcolsep) * \real{0.3333}}@{}}
\toprule\noalign{}
\begin{minipage}[b]{\linewidth}\raggedright
Item in this draft
\end{minipage} & \begin{minipage}[b]{\linewidth}\raggedright
Mathematical dependencies
\end{minipage} & \begin{minipage}[b]{\linewidth}\raggedright
Pinned source lineage
\end{minipage} \\
\midrule\noalign{}
\endhead
\bottomrule\noalign{}
\endlastfoot
Theorem O and cycle example & Ordered half selection, arc-product
factorization & ORDERED, ORDERED\_REVIEW, R08 \\
Corollary O and Lemma S & Flatness, unitary compression, singular-value
rank/energy & ORDERED; CURRENT, Section 2B; R07--R08 \\
Lemma N & Rank-nullity and an invertible full matrix & ROUTED\_PROOF,
CURRENT; classical nullity reference \\
Lemma C & Marked half-separator, Lemma N, two-color grouping & CURRENT,
Section 2A; CURRENT\_REVIEW; R07 \\
Lemma R & Per-row support union and integer division & CORRECTION; R08,
C05 \\
Theorem C and (19) & Lemmas C, R, S; legal recurrence and threshold
algebra & CORRECTION, CORRECTION\_RECEIPT; R08 \\
Lemma P and \(q=2\) & Column independence; pseudoforest structure;
weighted centroid & FIELD\_EXTENSION; CURRENT, Section 2A;
CURRENT\_REVIEW \\
Lemma F and Theorem F & Column deletion nullspace; endpoint rank
classes; Vandermonde & FIELD\_BASE, FIELD\_EXTENSION; R08 \\
Scaling transfer & Invertible diagonal normalization and norm
inequalities & CURRENT, Section 2B \\
Literature comparison & Bounded matched-statement inspection, not a
proof dependency & ROUTED\_PRIORITY and the linked primary texts \\
Current source-identity qualification & Recomputed source bytes and
recovered-body identity record & READINESS, SOURCE\_INDEX,
CORRECTION\_IDENTITY \\
\end{longtable}

The source named ROUTED\_PROOF is the older sealed \(36^r\)
presentation. It proves the same general structural approach with a less
efficient color selection. The current \(9^r\) and \(q=2\) proofs are in
CURRENT, independently reconstructed in R07 and discussed in R08. The
theorem body CORRECTION gives the combined sparse/rank/diagonal formula.
This draft re-exposes those existing results; it does not substitute the
old bound for the current one or infer a missing strengthening from a
short receipt.

CORRECTION\_RECEIPT is a declared attribution-normalized export. Its
recovered bytes have SHA-256
\nolinkurl{ac60fadc696be34bdea08f521b59714b2137d1a4e3874fe763578cfa062d0dfd}.
The recorded original canonical receipt identity is
\nolinkurl{eba23a76dd87fef48f9c5006bb416ee1756745874de212d83099ab1bc6af2b74};
those original receipt bytes are not present in this packet and are not
reconstructed. The correction theorem body itself is canonical-byte
matched. CORRECTION\_IDENTITY closes provenance only and explicitly does
not supply a new mathematical acceptance.

The companion SOURCE\_MANIFEST.json supplies each source's full SHA-256,
size, original local locator, packet copy, and role. All source copies
are inert text or JSON. Mixed-topic originals are included for faithful
provenance; only the sections identified by the dependency table support
this manuscript. Unrelated theorems and fusion sections in those
originals are outside its scope. No inherited research program was
executed for this consolidation. Fresh author-side arithmetic and
integrity checks, if present in the packet, are supplementary controls,
not an independent proof receipt.

\section{Frozen scope and remaining review
questions}\label{frozen-scope-and-remaining-review-questions}

No unresolved proof gap was found while reconstructing the included
statements. The final independent reviewer should still verify the exact
endpoint orientation, the marked-separator quantifier, the two-color
integer selection, the threshold legality, the component-nullspace
count, and the target normalization directly from this manuscript.

The following remain expressly open or excluded: exact historical
priority of the full contract, optimal constants and exponents, products
or interleavings of sums, general shared evaluators, arbitrary
arithmetic runtime, fixed-positive-tolerance logarithmic lower bounds,
arbitrary nonunitary target normalization without scale factors, and
extensions that change the additive grammar. No such problem is silently
assumed solved by consolidating the accepted proofs.

\end{document}
